\documentclass[11pt,a4paper]{article}

% Required packages
\usepackage[utf8]{inputenc}
\usepackage[T1]{fontenc}
\usepackage{amsmath}
\usepackage{graphicx}
\usepackage{booktabs}
\usepackage{hyperref}
\usepackage[margin=1in]{geometry}
\usepackage{natbib}
\usepackage{xcolor}
\usepackage{lineno}
\usepackage{authblk}

\bibliographystyle{unsrtnat}

\title{Caspase-11 Catalytic Activity and Autoprocessing Are Required for Noncanonical Inflammasome Speck Formation}

\author[1]{Author A}
\author[1]{Author B}
\author[1,2]{Author C}
\affil[1]{Department of Immunology, Harvard Medical School, Boston, MA, USA}
\affil[2]{Corresponding author: \texttt{author.c@hms.harvard.edu}}

\date{}

\begin{document}
\maketitle

\begin{abstract}
Caspase-11 (Casp11) is the central effector of the noncanonical inflammasome, sensing cytosolic lipopolysaccharide (LPS) through its CARD domain to trigger pyroptotic cell death. While Casp11 is known to oligomerize into a high-molecular-weight complex, the direct visualization of Casp11 supramolecular organizing centers (SMOCs/specks) and the molecular determinants governing their assembly have not been established. Here we develop a Casp11-mCherry fluorescent reporter and demonstrate that LPS transfection induces dose-dependent perinuclear speck formation in \textit{Casp11$^{-/-}$} bone marrow-derived macrophages reconstituted with wild-type Casp11-mCherry. The catalytically dead C254A mutant completely fails to form specks despite comparable expression, proving that CARD--LPS binding alone is insufficient for higher-order assembly. In HEK293T cells, we demonstrate that both catalytic activity (C254) and interdomain linker (IDL) cleavage at D285 are independently required for speck formation. Using an engineered TEV-cleavable Casp11, we show that exogenous protease cleavage at the IDL rescues speck formation even in catalytically dead mutants, confirming that the physical act of cleavage---not ongoing catalytic activity---is the key activating event. Co-transfection experiments further reveal that wild-type Casp11 rescues C254A-mCherry speck formation through trans-cleavage at D285, but this rescue requires the D285 site on the tagged molecule itself, establishing cis-autoprocessing as essential for oligomerization competence. These findings support a revised model in which LPS-induced dimerization activates catalytic activity, catalytic activity mediates intra-molecular autoprocessing at the IDL, and the processed form then assembles into the higher-order oligomeric SMOC.
\end{abstract}

\section{Introduction}

The noncanonical inflammasome pathway is a critical arm of innate immunity that detects cytosolic lipopolysaccharide (LPS) from Gram-negative bacteria independently of Toll-like receptor signaling \citep{Shi2014, Kayagaki2011, Hagar2013}. In mice, caspase-11 (Casp11) serves as both the sensor and the effector of this pathway: its N-terminal CARD domain directly binds the lipid A moiety of LPS, while its C-terminal catalytic domain cleaves gasdermin D (GSDMD) to trigger pyroptotic cell death \citep{Shi2015, Kayagaki2015, Aglietti2016}.

The activation of inflammatory caspases involves a proximity-induced dimerization model, in which recruitment to a signaling platform promotes dimerization and subsequent autoprocessing \citep{Salvesen2009, Boucher2018}. For the canonical inflammasome, caspase-1 is recruited to the ASC speck through CARD--CARD interactions, forming a large supramolecular organizing center (SMOC) that can be directly visualized as a perinuclear punctum by fluorescence microscopy \citep{Lu2014, Dick2016, Stutz2017}. In contrast, the assembly of the noncanonical Casp11 inflammasome has been inferred primarily from biochemical evidence---gel filtration showing high-molecular-weight complexes and Blue Native PAGE---but direct visualization of Casp11 specks in physiologically relevant cell types has been lacking \citep{Lee2018, Ross2018}.

Several molecular events are required for full Casp11 activation. LPS binding through the CARD domain is necessary but not sufficient: recent work demonstrated that both catalytic activity at C254 and autoprocessing at the IDL site D285 are required for GSDMD cleavage \citep{Lee2018, Ross2018, Broz2020}. However, whether these same molecular determinants control the upstream assembly of the inflammasome complex itself---as opposed to its downstream substrate processing activity---has remained unresolved. Furthermore, whether autoprocessing occurs in cis (intra-molecular) or in trans (inter-molecular) has important implications for the mechanism of inflammasome nucleation.

Here we develop a fluorescent Casp11-mCherry reporter system and use a combination of macrophage reconstitution, ectopic expression, structure-guided mutagenesis, and orthogonal protease rescue experiments to dissect the molecular requirements for noncanonical inflammasome speck formation.

\section{Results}

\subsection{LPS induces dose-dependent Casp11 speck formation in macrophages}

To directly visualize Casp11 SMOC formation, we generated a Casp11-mCherry fusion construct and reconstituted \textit{Casp11$^{-/-}$} bone marrow-derived macrophages (BMDMs) via lentiviral transduction (Table~\ref{tab:constructs}). Following priming with Pam3CSK4 and transfection with LPS from \textit{S. enterica} serovar Minnesota, we assessed pyroptosis by LDH release, GSDMD cleavage by western blot, and speck formation by confocal microscopy.

\begin{table}[htbp]
\centering
\caption{Constructs generated for this study.}
\label{tab:constructs}
\small
\begin{tabular}{lllll}
\toprule
\textbf{Construct} & \textbf{Tag} & \textbf{Backbone} & \textbf{Mutation} & \textbf{Purpose} \\
\midrule
Casp11 WT-mCherry & C-term mCherry & Lenti-SFFV & None & WT reporter \\
Casp11 C254A-mCherry & C-term mCherry & Lenti-SFFV & C254A & Catalytic dead \\
Casp11 D285A-mCherry & C-term mCherry & pReceiver-M56 & D285A & Non-cleavable \\
Casp11 C254A/D285A-mCherry & C-term mCherry & pReceiver-M56 & C254A + D285A & Double mutant \\
TEV-D285A-mCherry & C-term mCherry & Lenti-SFFV & D285A + TEV site & TEV rescue \\
TEV-C254A/D285A-mCherry & C-term mCherry & Lenti-SFFV & Both + TEV site & Neg. control \\
Untagged Casp11 WT & None & pReceiver-M56 & None & Trans donor \\
Untagged Casp11 C254A & None & pReceiver-M56 & C254A & Dead donor \\
Untagged Casp11 D285A & None & pReceiver-M56 & D285A & Active donor \\
\bottomrule
\end{tabular}
\end{table}

Wild-type Casp11-mCherry restored dose-dependent pyroptosis in \textit{Casp11$^{-/-}$} BMDMs, with LDH release increasing from approximately 20\% at 0.5~$\mu$g/mL LPS to approximately 65\% at 5~$\mu$g/mL (Table~\ref{tab:ldh}). The C254A-mCherry mutant showed no cytotoxicity at any dose. GSDMD cleavage was observed only in WT-mCherry reconstituted cells.

\begin{table}[htbp]
\centering
\caption{Cytotoxicity by LDH release in reconstituted \textit{Casp11$^{-/-}$} BMDMs.}
\label{tab:ldh}
\begin{tabular}{llll}
\toprule
\textbf{LPS ($\mu$g/mL)} & \textbf{WT-mCherry (\%)} & \textbf{C254A-mCherry (\%)} & \textbf{Empty Vector (\%)} \\
\midrule
0 & $\sim$5 & $\sim$5 & $\sim$5 \\
0.5 & $\sim$20 & $\sim$5 & $\sim$5 \\
1.0 & $\sim$35 & $\sim$5 & $\sim$5 \\
2.0 & $\sim$50 & $\sim$5 & $\sim$5 \\
5.0 & $\sim$65 & $\sim$5 & $\sim$5 \\
\bottomrule
\end{tabular}
\end{table}

Critically, confocal microscopy revealed that WT-mCherry cells formed large, bright perinuclear specks after LPS transfection in a dose-dependent manner, while C254A-mCherry remained entirely diffuse throughout the cytoplasm (Table~\ref{tab:bmdm_specks}).

\begin{table}[htbp]
\centering
\caption{Speck formation in LPS-transfected BMDMs.}
\label{tab:bmdm_specks}
\begin{tabular}{lll}
\toprule
\textbf{LPS ($\mu$g/mL)} & \textbf{WT-mCherry Specks (\%)} & \textbf{C254A-mCherry Specks (\%)} \\
\midrule
0 & $\sim$0 & $\sim$0 \\
0.5 & $\sim$8 & $\sim$0 \\
1.0 & $\sim$15 & $\sim$0 \\
2.0 & $\sim$25 & $\sim$0 \\
5.0 & $\sim$35 & $\sim$0 \\
\bottomrule
\end{tabular}
\end{table}

This demonstrates that CARD--LPS binding alone is insufficient for speck formation; catalytic activity at C254 is essential for higher-order SMOC assembly.

\subsection{Catalytic activity is required for spontaneous speck formation in HEK293T cells}

To facilitate mechanistic dissection with mutant constructs, we established an ectopic expression system in HEK293T cells. Overexpression of WT-mCherry produced spontaneous specks in approximately 40--50\% of mCherry-positive cells, while C254A-mCherry produced specks in only 2--5\% (Table~\ref{tab:hek_specks}).

\begin{table}[htbp]
\centering
\caption{Spontaneous speck formation in HEK293T cells.}
\label{tab:hek_specks}
\begin{tabular}{lll}
\toprule
\textbf{Construct} & \textbf{Specks (\%)} & \textbf{$N$} \\
\midrule
WT-mCherry & $\sim$40--50 & 3 \\
C254A-mCherry & $\sim$2--5 & 3 \\
\bottomrule
\end{tabular}
\end{table}

Blue Native PAGE confirmed that WT-mCherry assembled into complexes exceeding 1~MDa, while C254A-mCherry remained as monomer/dimer species \citep{Boucher2018}.

\subsection{IDL cleavage at D285 is independently required for speck formation}

To determine whether autoprocessing at the IDL is required for oligomerization, we tested the D285A non-cleavable mutant. D285A-mCherry phenocopied the catalytic-dead C254A mutant, forming specks in only 2--5\% of cells (Table~\ref{tab:d285a}). The C254A/D285A double mutant also failed to form specks, confirming that both catalytic activity and physical cleavage at D285 are independently necessary.

\begin{table}[htbp]
\centering
\caption{IDL cleavage requirement for speck formation.}
\label{tab:d285a}
\begin{tabular}{ll}
\toprule
\textbf{Construct} & \textbf{Specks (\%)} \\
\midrule
WT-mCherry & $\sim$40--50 \\
D285A-mCherry & $\sim$2--5 \\
C254A/D285A-mCherry & $\sim$2--5 \\
\bottomrule
\end{tabular}
\end{table}

\subsection{Wild-type Casp11 rescues C254A speck formation through trans-cleavage}

To test whether trans-processing can rescue speck formation, we co-transfected fixed amounts of C254A-mCherry with increasing doses of untagged wild-type Casp11. Specks were scored only in the mCherry channel to track the catalytically dead reporter (Table~\ref{tab:trans_rescue}).

\begin{table}[htbp]
\centering
\caption{Trans-rescue of C254A-mCherry by co-expression with untagged WT Casp11.}
\label{tab:trans_rescue}
\begin{tabular}{lll}
\toprule
\textbf{Untagged WT Dose} & \textbf{C254A-mCherry Specks (\%)} & \textbf{WT-mCherry Control (\%)} \\
\midrule
0 & $\sim$2 & $\sim$45 \\
0.25$\times$ & $\sim$10 & $\sim$45 \\
0.5$\times$ & $\sim$18 & $\sim$48 \\
1.0$\times$ & $\sim$25 & $\sim$50 \\
2.0$\times$ & $\sim$30 & $\sim$50 \\
\bottomrule
\end{tabular}
\end{table}

Critically, the C254A/D285A double mutant was not rescued by co-expression with wild-type Casp11 (Table~\ref{tab:double_rescue}), demonstrating that the rescue requires cleavage at D285 on the tagged molecule itself.

\begin{table}[htbp]
\centering
\caption{The C254A/D285A double mutant cannot be rescued by WT Casp11.}
\label{tab:double_rescue}
\begin{tabular}{ll}
\toprule
\textbf{Untagged WT Dose} & \textbf{C254A/D285A-mCherry Specks (\%)} \\
\midrule
0 & $\sim$2 \\
0.25$\times$ & $\sim$2 \\
0.5$\times$ & $\sim$3 \\
1.0$\times$ & $\sim$3 \\
2.0$\times$ & $\sim$3 \\
\bottomrule
\end{tabular}
\end{table}

\subsection{Cis versus trans processing dissection}

To further distinguish cis from trans processing, we co-transfected C254A-mCherry with untagged WT, C254A, or D285A Casp11 (Table~\ref{tab:cis_trans}).

\begin{table}[htbp]
\centering
\caption{Cis versus trans processing: effect of different untagged co-transfectants.}
\label{tab:cis_trans}
\begin{tabular}{lll}
\toprule
\textbf{Untagged Donor} & \textbf{C254A-mCherry Processing?} & \textbf{Speck Rescue?} \\
\midrule
WT Casp11 & Yes & Yes (dose-dependent) \\
C254A Casp11 & No & No \\
D285A Casp11 & Yes & Yes (dose-dependent) \\
\bottomrule
\end{tabular}
\end{table}

The D285A untagged donor---which is catalytically active but itself non-cleavable---rescued both processing and speck formation of C254A-mCherry. This proves that the catalytic activity of the untagged partner cleaves C254A-mCherry in trans at D285, and the resulting processed C254A-mCherry then assembles through cis-dependent conformational changes.

\subsection{TEV protease cleavage rescues speck formation in catalytically dead mutants}

The key experiment distinguishing ``catalytic activity needed continuously'' from ``cleavage is the activating event'' used an engineered TEV-cleavable Casp11. A TEV protease cleavage site (ENLYFQ/G) was inserted at the IDL of the D285A mutant (Table~\ref{tab:tev}).

\begin{table}[htbp]
\centering
\caption{TEV protease rescue of speck formation.}
\label{tab:tev}
\begin{tabular}{llll}
\toprule
\textbf{Construct} & \textbf{TEV Protease} & \textbf{Processing?} & \textbf{Specks?} \\
\midrule
TEV-D285A-mCherry & $-$ & No & No \\
TEV-D285A-mCherry & $+$ & Yes & Yes ($\sim$30--40\%) \\
TEV-C254A/D285A-mCherry & $-$ & No & No \\
TEV-C254A/D285A-mCherry & $+$ & Yes & Yes ($\sim$25--35\%) \\
\bottomrule
\end{tabular}
\end{table}

Even the catalytically dead TEV-C254A/D285A construct was rescued by TEV protease, confirming that it is the physical act of IDL cleavage---not ongoing Casp11 catalytic activity---that enables higher-order oligomerization.

\subsection{Comprehensive summary of speck formation determinants}

Table~\ref{tab:master_summary} integrates all experimental findings across all constructs and assay systems.

\begin{table}[htbp]
\centering
\caption{Summary of speck formation across all constructs and conditions.}
\label{tab:master_summary}
\small
\begin{tabular}{lllllll}
\toprule
\textbf{Construct} & \textbf{Cat.} & \textbf{IDL} & \textbf{BMDMs} & \textbf{HEK293T} & \textbf{WT Rescue} & \textbf{TEV Rescue} \\
\midrule
WT & Yes & Yes & Yes & Yes & N/A & N/A \\
C254A & No & Yes & No & No & Yes & N/A \\
D285A & Yes & No & N/D & No & N/D & N/A \\
C254A/D285A & No & No & N/D & No & No & N/A \\
TEV-D285A & Yes & TEV & N/D & No/Yes & N/D & Yes \\
TEV-C254A/D285A & No & TEV & N/D & No/Yes & N/D & Yes \\
\bottomrule
\end{tabular}
\end{table}

\section{Discussion}

Our study provides the first direct visualization of Casp11 SMOC formation in primary macrophages and establishes the molecular hierarchy governing noncanonical inflammasome assembly. The revised mechanistic model (Table~\ref{tab:model}) proceeds through a series of biochemically separable steps: LPS binding to the CARD domain, dimerization, catalytic activation, intra-molecular autoprocessing at the IDL, conformational change, and higher-order oligomerization into the visible SMOC.

\begin{table}[htbp]
\centering
\caption{Proposed mechanistic model for noncanonical inflammasome assembly.}
\label{tab:model}
\begin{tabular}{lll}
\toprule
\textbf{Step} & \textbf{Event} & \textbf{Evidence} \\
\midrule
1 & LPS binds Casp11 CARD & Published \citep{Shi2014} \\
2 & LPS-bound Casp11 dimerizes & Published \citep{Ross2018} \\
3 & Dimerization activates catalysis & Published \\
4 & Autoprocessing at D285 (IDL) & This study: C254A/D285A block \\
5 & Conformational change & This study: TEV rescue \\
6 & Oligomeric SMOC forms & This study: microscopy \\
7 & SMOC cleaves GSDMD & Published + this study \\
8 & GSDMD-NT pores / pyroptosis & Published \\
\bottomrule
\end{tabular}
\end{table}

The most informative experiments are the TEV protease rescue of catalytically dead mutants and the D285A untagged donor rescue of C254A-mCherry. Together, these demonstrate that: (1) the physical act of cleavage at the IDL, not ongoing catalytic activity, is the activating event for oligomerization; and (2) trans-cleavage at D285 can occur between Casp11 molecules, but the resulting processed molecule must undergo cis-dependent conformational changes to become oligomerization-competent.

This model has important implications for the kinetics of noncanonical inflammasome assembly. The requirement for autoprocessing before oligomerization creates a kinetic barrier that may serve as a regulatory checkpoint, preventing spurious speck formation from transient LPS encounters \citep{Boucher2018, Broz2020}. The dose-dependent nature of speck formation in BMDMs is consistent with this threshold model: low LPS concentrations may be sufficient for some dimerization and catalytic activation but insufficient to drive the autoprocessing required for SMOC assembly.

Our findings also have broader implications for understanding inflammatory caspase activation mechanisms. The canonical caspase-1 pathway involves recruitment to an ASC speck that serves as the organizing platform \citep{Lu2014, Dick2016}. In contrast, Casp11 must self-organize into its own SMOC through an intrinsic autoprocessing-dependent mechanism. This fundamental difference may explain why the noncanonical pathway exhibits different activation kinetics and threshold behavior compared to the canonical pathway \citep{Kayagaki2011, Hagar2013}.

\section{Materials and Methods}

\subsection{Cell culture and differentiation}
BMDMs were differentiated from \textit{Casp11$^{-/-}$} (C57BL/6 background) bone marrow in 30\% L929-conditioned DMEM for 7 days \citep{Kayagaki2011}. HEK293T cells were maintained in standard DMEM with 10\% FBS (Table~\ref{tab:conditions}).

\begin{table}[htbp]
\centering
\caption{Experimental conditions and reagents.}
\label{tab:conditions}
\begin{tabular}{ll}
\toprule
\textbf{Parameter} & \textbf{Detail} \\
\midrule
BMDM differentiation & 30\% L929-conditioned DMEM, 7 days \\
Priming agent & Pam3CSK4, 4 h \\
LPS source & \textit{S. enterica} serovar Minnesota \\
LPS delivery & Lipofectamine/FuGENE transfection \\
Microscopy fixation & 6 h post-LPS \\
LDH assay & 16 h post-LPS \\
HEK293T imaging & 18 h post-transfection \\
Transfection (HEK293T) & PEI, 1:1 w/w DNA:PEI \\
Lentiviral packaging & pVSV-G + psPAX2 \\
Microscopy & Confocal, Hoechst nuclear stain \\
Western blot targets & mCherry, GSDMD, $\beta$-actin \\
Mutagenesis & Q5 SDM Kit (NEB \#E0554S) \\
TEV consensus & ENLYFQ/G \\
\bottomrule
\end{tabular}
\end{table}

\subsection{Lentiviral transduction}
Lentiviral particles were produced in HEK293T cells using pTwist Lenti-SFFV-WPRE vectors co-transfected with pVSV-G and psPAX2 packaging plasmids. BMDMs were transduced at day 2 of differentiation and used at day 7.

\subsection{Speck quantification}
Cells were fixed with 4\% paraformaldehyde, stained with Hoechst, and imaged by confocal fluorescence microscopy. Specks were defined as bright perinuclear mCherry puncta. The percentage of mCherry-positive cells containing at least one speck was calculated from at least 100 mCherry-positive cells per condition.

\subsection{Western blotting}
Cell lysates were resolved by SDS-PAGE (reducing and non-reducing) and Blue Native PAGE. Blots were probed with anti-mCherry, anti-GSDMD, and anti-$\beta$-actin antibodies.

\subsection{TEV protease rescue experiments}
A TEV protease cleavage site (ENLYFQ/G) was introduced at the IDL of D285A-mCherry and C254A/D285A-mCherry constructs using the Q5 Site-Directed Mutagenesis Kit. Constructs were co-transfected with a TEV protease expression plasmid in HEK293T cells. Processing and speck formation were assessed at 18~h post-transfection.

\subsection{Co-transfection experiments}
Fixed amounts (300~ng) of tagged mCherry reporters were co-transfected with varying doses of untagged Casp11 constructs (0, 75, 150, 300, 600~ng) using PEI at 1:1 w/w DNA:PEI ratio. Total DNA was equalized with empty vector. Specks were scored exclusively in the mCherry channel.

\section*{Acknowledgments}
We thank members of the laboratory for helpful discussions and the Harvard Medical School Imaging Core for microscopy support.

\bibliography{references}

\end{document}
